\chapter{Testing}

\section{What is worth testing}

The project has around 118 tests, and they are not spread evenly. They
concentrate on three kinds of code:

\begin{enumerate}
  \item \textbf{Arithmetic that decides money.} Profit per broker, percentages,
        the paste parser, aggregation over ranges.
  \item \textbf{Parsing of external formats.} Broker files that the project does
        not control and cannot fix.
  \item \textbf{Contracts that are easy to break by accident.} The health
        endpoints, the task token, what the backup must never contain.
\end{enumerate}

There are no tests asserting that a template contains a particular heading.
Those break on every wording change and catch nothing.

\section{The database fixture}

Most tests run against SQLite in memory, created and destroyed per test:

\begin{lstlisting}[style=py]
@pytest.fixture
def db():
    engine = create_engine("sqlite:///:memory:", connect_args={"check_same_thread": False})
    Base.metadata.create_all(engine)
    s = sessionmaker(bind=engine)()
    try:
        yield s
    finally:
        s.close()
\end{lstlisting}

Tests that go through \code{TestClient} need one extra detail:

\begin{lstlisting}[style=py]
engine = create_engine("sqlite:///:memory:",
                       connect_args={"check_same_thread": False},
                       poolclass=StaticPool)
\end{lstlisting}

\code{StaticPool} forces every connection to be the same one. Without it, the
test client runs the request in another thread, gets a fresh connection, and an
in-memory database that is empty, because each connection to
\code{:memory:} is a different database.

\begin{keypoint}[A test failing for the right reason]
An early version of the security test used the developer's local \code{dev.db}
by accident. It passed on that machine and failed on a clean clone with
\code{no such table: users}. The fixture was rewritten to inject an in-memory
database through \code{dependency\_overrides}, and a comment in the test records
why.
\end{keypoint}

\section{Overriding dependencies}

FastAPI dependency injection makes it easy to replace the database and the
authenticated user:

\begin{lstlisting}[style=py]
app.dependency_overrides[get_session] = lambda: db
app.dependency_overrides[require_user] = lambda: user
monkeypatch.setattr(inv_router, "validate_csrf", lambda *a, **k: True)
\end{lstlisting}

The CSRF check is patched because these tests exercise the handler logic, not
the CSRF mechanism, which has its own dedicated tests. Being explicit about that
split keeps each test focused on one thing.

\section{No network in the test suite}

Three parts of the application talk to the internet: prices, ISIN lookup and
Telegram. All three are patched at the boundary.

\begin{lstlisting}[style=py]
monkeypatch.setattr(prices, "fetch_price", lambda s: Decimal("50000"))
\end{lstlisting}

\begin{lstlisting}[style=py]
async def fake_send(chat_id, text, reply_markup=None):
    msgs.append((chat_id, text, reply_markup))
monkeypatch.setattr(bot, "send_message", fake_send)
\end{lstlisting}

Capturing the outgoing messages in a list turns "the bot replied correctly" into
a plain assertion on strings, and the whole suite runs offline in seconds.

\section{Parser tests without private data}

The private version of the project tested the parsers against real statements
stored outside the repository, and skipped when they were absent. That is
convenient and useless to anyone else: on any other machine the tests silently
skip.

The public version builds small synthetic statements inside the test file:

\begin{lstlisting}[style=py]
KRAKEN_CSV = b"""txid,pair,time,type,ordertype,price,cost,fee,vol
1,BTC/EUR,2026-01-02,buy,market,40000,400.00,0.4,0.01
2,BTC/EUR,2026-02-02,buy,market,45000,450.00,0.4,0.01
3,BTC/EUR,2026-03-02,sell,market,50000,250.00,0.2,0.005
"""

def test_kraken_nets_buys_and_sells():
    r = kraken.parse(KRAKEN_CSV)
    btc = next(p for p in r.positions if p.asset == "BTC/EUR")
    assert btc.invested == Decimal("600.00")      # 400 + 450 - 250
    assert btc.quantity == Decimal("0.01500000")
\end{lstlisting}

The data is fake, the format is real, and the test runs everywhere. It also
documents the format better than a paragraph would.

\section{The startup test}

One test does not check behaviour, it protects a property that is easy to lose:

\begin{lstlisting}[style=py]
HEAVY = ["matplotlib", "reportlab", "openpyxl", "pandas", "pdfplumber", "yfinance"]

PROBE = f"""
import json, sys
sys.path.insert(0, {str(ROOT)!r})
import app.main
print(json.dumps([m for m in {HEAVY!r} if m in sys.modules]))
"""

def test_boot_does_not_import_heavy_libraries():
    out = subprocess.run([sys.executable, "-c", PROBE], capture_output=True,
                         text=True, cwd=ROOT, check=True)
    loaded = json.loads(out.stdout.strip().splitlines()[-1])
    assert loaded == [], (
        f"Importing app.main loads {loaded}. Move those imports inside the "
        "function that uses them (see app/routers/reports.py).")
\end{lstlisting}

It runs in a clean subprocess because inside pytest those libraries are already
loaded by other tests, so checking \code{sys.modules} in-process would always
fail.

The failure message is as important as the assertion: it names the offending
modules and points at the fix. A future contributor adding a top-level
\code{import pandas} gets told exactly what to do instead of wondering why an
unrelated test broke.

\section{Regression tests carry their reason}

When a bug is fixed, the test that pins it keeps the explanation:

\begin{lstlisting}[style=py]
def test_donut_drops_negative_slices():
    """A matplotlib pie raises on slices <= 0; donut has to filter them out."""
\end{lstlisting}

Six months later, that docstring is the difference between deleting an odd test
and understanding why it exists.

\section{Continuous integration}

A GitHub Actions workflow runs the suite on every push and pull request:

\begin{lstlisting}[style=sh]
- uses: actions/setup-python@v5
  with:
    python-version: "3.12"
    cache: pip
- run: pip install -r requirements.txt
- run: python -m pytest -q
\end{lstlisting}

There is no coverage gate and no linter in the pipeline. The suite either passes
or it does not, which for a project of this size is the check that carries the
information.
